\section{Discussion}\label{discussion}

The preceding six acts established that LLM conceptual navigation has measurable geometric structure (Section 4), consistent with quasimetric topology (Section 5), causally sensitive to early anchoring (Section 6), resistant to single-variable mechanistic explanation (Section 7), generalizable across 12 models from 11 independent training pipelines (Section 8), and robust to protocol variation (Section 9). This section synthesizes those findings, states their limitations directly, and identifies the directions they open. Across approximately 21,540 runs, 12 models, 11 phases, and 29 documented dead ends, the benchmark reveals a consistent empirical picture: conceptual navigation is structured, asymmetric, model-specific in dynamics, and opaque in mechanism.

\subsection{Summary of the Six-Act Narrative}\label{summary-of-the-six-act-narrative}

\textbf{Structure (Section 4).} Models have distinct, stable conceptual gaits spanning a 2.9x range from GPT (0.258) to Mistral (0.747) {[}robust{]} R1. Navigation is self-consistent across time and resolution --- cross-batch Jaccard falls within 0.05 of within-batch (O9), and 5-waypoint paths appear as proper subsequences of 10-waypoint paths 67.9\% of the time (O10). The dual-anchor effect confirms that both endpoints shape the trajectory, ruling out pure forward-chaining (O2).

\textbf{Topology (Section 5).} Conceptual space is consistent with quasimetric structure. The symmetry axiom fails comprehensively --- mean directional asymmetry is 0.811 in the Phase 2 cohort, exceeding 0.60 for all 12 models tested {[}robust{]} R2. The triangle inequality holds at approximately 91\% across three independent samples, tested on within-direction Jaccard distance. These two properties are measured on different distance operationalizations (cross-direction Jaccard for asymmetry, within-direction Jaccard for triangle inequality), so the quasimetric characterization is convergent rather than formally established on a single unified metric. Bridge concepts function as structural bottlenecks --- obligatory intermediaries rather than mere associations {[}robust{]} R6. Hierarchical paths are compositional, with 4.9x higher transitivity than random controls {[}robust{]} R4.

\textbf{Mechanism (Section 6).} Pre-filling the first waypoint causally displaces bridges (O15, mean displacement 0.515, CI {[}0.357, 0.664{]} excluding zero). The primary mechanism is associative primacy: whatever occupies position 1 anchors the trajectory regardless of its semantic content, with displacement magnitudes statistically indistinguishable across incongruent, congruent, and neutral conditions. Relation class matters as a secondary modulator (O26): related pre-fills (on-axis or same-domain) preserve the navigational context and allow bridge recovery, while unrelated pre-fills destroy it (Friedman p = 0.034). Bridge survival under perturbation is bimodal --- ``sadness'' and ``dog'' survive robustly; ``harmony'' and ``germination'' collapse to near-zero.

\textbf{Limits (Section 7).} Seven follow-up hypotheses were falsified --- spanning mechanistic, model-deficit, methodological, and effect-direction types; one initially falsified result was resurrected with additional data. Prediction accuracy descended from 81.3\% in the characterization phases to 20-24\% in the mechanism phases (O24). The pattern is systematic: qualitative structure is characterizable and replicable at approximately 80\%; quantitative mechanism is opaque, with single-variable models succeeding at approximately 0-15\%.

\textbf{Generality (Section 8).} Gait and asymmetry replicate universally across 12 models from 11 independent training pipelines {[}robust{]} R1, R2. Bridge frequency generalizes among frontier-scale models (aggregate difference CI includes zero) but not to Llama 3.1 8B, which navigates through different landmarks (bridge frequency 0.200 vs 0.717-0.817 for frontier models). The within-family Llama comparison is particularly informative: Maverick (frontier-scale) shows bridge frequency 0.724 while 8B shows 0.200, a 3.6x difference. The hierarchy is structure \textgreater{} content \textgreater{} scale: structural properties are universal, navigational content is shared among large models, and scale is the clearest differentiator {[}robust{]} R6 extended.

\textbf{Robustness (Section 9).} Model identity dominates protocol variation in an ANOVA decomposition (eta-squared = 0.242 for model vs approximately 0 for waypoint count and temperature) (O30). Bridge frequency is the most robust property, exceeding 0.97 across all conditions including the most challenging (5 waypoints, temperature 0.9) (O31). Gait rank ordering is largely stable (Kendall's W = 0.840), with Claude maintaining the top position in every condition. Asymmetry is resolution-dependent: 9-waypoint conditions clear the 0.60 threshold (0.669-0.684) while 5-waypoint conditions fall just short (0.593-0.594) (O32). Temperature effects on all metrics are negligible.

\subsection{What the Benchmark Reveals About Conceptual Space}\label{what-the-benchmark-reveals-about-conceptual-space}

Four interpretive points emerge from the empirical results.

\textbf{Consistent with quasimetric structure.} Symmetry fails comprehensively (mean 0.811). The triangle inequality holds at approximately 91\% across three independent samples. These two findings use different operationalizations of distance --- cross-direction Jaccard for asymmetry, within-direction Jaccard for triangle inequality --- so the formal quasimetric claim is ``consistent with'' rather than deductively established on a single metric. The asymmetry finding parallels two independent results from different literatures: \citeauthor{tversky1977features}'s (\citeyear{tversky1977features}) demonstration that human similarity judgments are inherently asymmetric (North Korea is judged more similar to China than China is to North Korea), and \citeauthor{berglund2024reversal}'s (\citeyear{berglund2024reversal}) reversal curse showing that autoregressive models cannot reverse learned associations. Our finding is consistent with both: navigational asymmetry parallels the reversal curse and Tversky's human similarity asymmetry. The connection is thematic --- all three concern directional asymmetry in conceptual structure --- rather than mechanistic; we do not demonstrate that the reversal curse causes the specific asymmetry we observe.

\textbf{Bottleneck topology.} The most reliable navigational landmarks are not the most strongly associated concepts but obligatory intermediaries: ``spectrum'' for light-to-color (frequency 1.00 across all four original models), ``dog'' for animal-to-poodle, ``warm'' for hot-to-cold. ``Metaphor'' --- intuitively central to the language-to-thought relationship --- appeared at 0.00 frequency (graveyard entry G5). ``Energy'' --- associated with both hot and cold --- appeared at 0.00 on hot-to-cold. Process-naming concepts outperformed object-naming concepts: ``germination'' outperformed ``plant'' for seed-to-garden in all four models. This echoes a finding from knowledge graph research: shortest-path intermediaries are more informative than high-degree hubs. Navigational salience depends on directional information content, not on generic associative strength.

\textbf{Model-specific gaits on shared geometry.} All frontier-scale models navigate the same basic geometry through the same strongest landmarks, but with different route preferences. Mistral (gait 0.747) produces near-deterministic paths --- 0.936 Jaccard on music-to-mathematics, producing nearly identical waypoint sets across all repetitions. GPT (0.258) explores broadly, sampling different landmarks on each traversal. The aggregate bridge frequency difference between original and expanded cohorts has a CI including zero, yet gait varies 2.9x. This is the behavioral analog of a refinement to the Platonic Representation Hypothesis: representations may converge \citep{huh2024platonic}, but navigation diversifies. Same map, different routes. Structure converges; gaits remain model-specific. Neither extreme is intrinsically superior --- they represent genuinely different navigational strategies that a model-agnostic benchmark must accommodate.

\textbf{Navigational structure is not compositional from simple variables.} The mechanism ceiling --- 7 of 8 hypotheses falsified at primary test --- suggests that navigation emerges from the interaction of multiple factors rather than any single structural variable. Competitor count (G20), dominance ratio (G23), gradient type (G21, G24), and gait normalization (G22) all failed individually. The prediction accuracy descent from 81.3\% to 20\% tracks a shift in prediction type, not a loss of signal: replication predictions succeed at approximately 80\% throughout, while mechanistic predictions hover near 0-15\%. This parallels findings in computational neuroscience, where spatial navigation in rodents and humans depends on multi-factorial interactions between grid cells, place cells, head-direction cells, and boundary cells. The analogy is structural, not mechanistic: in both cases, the emergent navigational behavior resists decomposition into single-variable explanations.

\textbf{Statistical caveat on the quasimetric characterization.} The two tests supporting the quasimetric claim --- asymmetry and triangle inequality --- use different distance operationalizations. Asymmetry uses cross-direction Jaccard (comparing A-to-B paths against B-to-A paths); triangle inequality uses within-direction Jaccard (consistency of repeated same-direction runs). Additionally, all metrics are lexical --- synonymous waypoints are treated as distinct. The magnitude of the resulting bias is bounded by the canonicalization pipeline but is otherwise unknown. These caveats do not undermine the qualitative finding (systematic asymmetry, triangle inequality holding at approximately 91\%), but the formal quasimetric claim should be understood as ``consistent with'' rather than ``formally established.''

\subsection{Limitations}\label{limitations}

The benchmark's limitations divide into two distinct categories: design issues with the benchmark itself, and findings about the phenomenon that constrain interpretation. Conflating these would obscure the distinction between things we should have done better and things that turned out to be harder than expected.

\subsubsection{Benchmark Limitations}\label{benchmark-limitations}

\textbf{1. Control pair failure {[}robust{]} R5 (qualified).} This is the benchmark's most serious design limitation. The original control pair stapler-monsoon fails strict criteria for all 12 models (top waypoint frequency 0.650-1.000). Phase 11B tested 4 replacement candidates (accordion-stalactite, turmeric-trigonometry, barnacle-sonnet, magnesium-ballet) across 6 models: 0 of 24 cells passed either the frequency gate (\textless{} 0.15) or the entropy gate (\textgreater{} 5.0). Each candidate revealed creative structured routing --- turmeric-trigonometry produced ``ancient India'' (Claude), ``spice'' (GPT), ``golden ratio'' (Mistral) --- yet each model still converged on a single dominant waypoint. LLMs find navigable routes between any concept pair at 7 waypoints, defeating single-pair strict-threshold control validation. The benchmark's validity rests on the predictability gap --- experimental bridges are predicted from the endpoint relationship; control bridges are ad hoc --- not on presence versus absence of structure. This limitation should be weighted prominently in any evaluation of the benchmark's claims.

\textbf{2. Self-grounding circularity.} Model outputs generate their own ground truth. The concern is that structural findings could be artifacts of the generation process rather than properties of conceptual organization. Three mitigations partially address this:

\begin{enumerate}
\def\labelenumi{(\alph{enumi})}
\tightlist
\item
  Testing against external mathematical properties that have model-independent definitions (metric axioms --- the triangle inequality either holds or it does not, independent of the generation process).
\item
  Cross-model comparison across 12 independent models (if all 12 show the same property, it cannot be attributed to idiosyncrasies of any single model's decoding).
\item
  Causal intervention via the pre-fill experiments (bridge displacement demonstrates that structure is causally sensitive, not merely a statistical regularity).
\end{enumerate}

These mitigations reduce the concern but never fully eliminate it. The benchmark cannot distinguish between ``models have navigational geometry'' and ``the generation process has statistical regularities that mimic navigational geometry'' --- though the causal sensitivity of bridge displacement (Section 6.1) makes the second interpretation less parsimonious.

\textbf{3. Lexical metrics only.} All metrics use Jaccard similarity after lemmatization. ``Melody'' and ``tune'' are treated as entirely distinct. The asymmetry mean of 0.811 is lexical asymmetry; true conceptual asymmetry is likely lower because some apparent non-overlap reflects synonymous waypoints scored as different. The magnitude of this bias is unknown but bounded by the canonicalization pipeline, which handles morphological variants (running/run, cities/city).

Embedding-based semantic similarity was implemented but deferred due to API cost. This is the benchmark's most pervasive methodological limitation, present in every metric and every phase. It affects all structural claims: gait underestimates true conceptual consistency, asymmetry overestimates true conceptual asymmetry, and bridge frequency may miss bridges expressed through synonyms rather than exact tokens.

\textbf{4. API-mediated access.} Internal model state cannot be controlled, temperature settings may be approximate for some providers, and silent model updates may occur during multi-day data collection. Partially addressed by the no-temporal-drift finding (O9, cross-batch Jaccard within 0.05 of within-batch) and the robustness to temperature variation (Phase 11C, eta-squared = 0.002 for temperature). These mitigations apply to the tested conditions; they do not rule out sensitivity to untested API-level changes.

\textbf{5. Concept pair selection bias.} Pairs were selected by the experimenters based on semantic interest and hypothesis-driven design, not sampled from a principled distribution over conceptual space. The benchmark tests navigational structure on the pairs it tests, not on a representative sample of all possible concept relationships. Partially addressed by the breadth of the pair inventory (100+ unique pairs across 11 phases) and the generality analysis (Section 8), but the selection is not random and the results should not be extrapolated to arbitrary concept pairs without further testing.

\textbf{6. Resolution dependence.} Asymmetry requires 9 or more waypoints to reliably exceed the 0.60 threshold in the robustness subset (O32): the 5-waypoint conditions fall just short (0.593-0.594), the 7-waypoint baseline also fell short (0.599) in the sparse 3-model subset, and only the 9-waypoint conditions clear it (0.669-0.684). The full Phase 2 cohort at 5 waypoints produced 0.811, but that used a larger pair set and model cohort. Other structural properties may have similar resolution thresholds that have not been identified --- the benchmark tested waypoint counts of 5, 7, and 9 but not the full range. This is a measurement sensitivity finding, not evidence against asymmetry, but it means that claims about asymmetry magnitude are conditioned on sufficient path length.

\subsubsection{Phenomenon Limitations}\label{phenomenon-limitations}

\textbf{7. Mechanism ceiling.} Seven of eight hypotheses were falsified at the primary-test level (Section 7.2). This is not a design flaw --- it is the benchmark's most informative negative result. The mechanism ceiling establishes that conceptual navigation resists single-variable explanation and that the resolving power of one-factor experimental designs has been exhausted for this phenomenon. The prediction accuracy trajectory --- 81.3\% (Phase 4) to 20\% (Phase 9) --- tracks a systematic shift in prediction type, not random degradation.

\textbf{8. Cross-model distance failure (O18, O19).} Models rank-order navigational distances differently. The cross-model distance correlation was r = 0.170, far below the 0.50 validity threshold. Gait normalization produced zero improvement (r = 0.212 before and after, with 0.000 change on every model pair, G22). Some model pairs actively anti-correlate: Grok-Gemini at r = -0.580. No universal distance function is recoverable from path-based measurements. Within-model distance analysis remains viable; cross-model distance comparison is blocked.

\textbf{9. Gemini characterization failure.} Gemini's approximately 30\% bridge frequency deficit relative to the non-Gemini mean (0.480 vs approximately 0.67) is real and stable across the tested phases. Three successive mechanistic characterizations --- frame-crossing threshold (G7), gradient blindness (G21), transformation-chain blindness (G24) --- all failed, with two producing results in the opposite direction from predicted. Whatever pair type is designated as ``hard'' for Gemini, Gemini does relatively better on it. The type-based approach to characterizing the deficit appears fundamentally misguided; the deficit is model-level, not type-level. Future approaches might require representational analysis rather than behavioral type-classification.

\subsection{The Graveyard as Methodology}\label{the-graveyard-as-methodology}

The benchmark documents 29 dead ends across 11 phases --- a transparent record of the benchmark's learning curve.

Four lessons emerge from the graveyard.

First, intuitive semantic importance has zero predictive power for navigational behavior. ``Metaphor'' --- the concept most intuitively connected to both language and thought --- appeared at 0.00 frequency as a bridge for language-to-thought (G5). ``Energy'' --- associated with both hot and cold --- appeared at 0.00 on hot-to-cold. Navigational salience and human intuition about conceptual centrality are orthogonal; bridges must name the primary axis of connection, not merely inhabit the shared associative neighborhood.

Second, qualitative models can be correct while their quantitative predictions fail. Gemini's deficit is real (approximately 30\% reduction in bridge frequency), but three quantitative characterizations of its mechanism failed (G7, G21, G24). The qualitative observation replicates across every phase in which Gemini was tested; the mechanism remains opaque. This pattern --- robust qualitative finding, intractable quantitative mechanism --- characterizes the benchmark as a whole.

Third, retrospective signals on small samples do not generalize. The dominance ratio showed a promising correlation (rho = 0.60) on the original 8-pair sample but collapsed to rho = 0.157 on an independent sample (G23). The competitor count predictor showed a similar failure (G20, rho = 0.116). Both variables were plausible on the original data; neither survived replication.

Fourth, infrastructure limitations can masquerade as substantive findings. The bridge generalization failure (G26) was initially attributed to the phenomenon based on data from a single small model (Llama 8B) under aggressive 60-second timeouts that excluded most new models on latency, not capability. After timeout relaxation to 300 seconds and testing of 3 additional frontier-scale models, the aggregate finding reversed --- bridge frequency CI shifted from excluding zero to including it. Llama 8B's low bridge frequency (0.200) remains real as a scale effect, but the original conclusion that bridge structure does not generalize was wrong. G26 is the sole clean resurrection in the graveyard and a caution against both premature interpretation of negative results and confounding infrastructure limitations with phenomenon limitations.

Transparent dead-end documentation should be standard practice in benchmarking work. The incentive structure in machine learning research favors reporting only successful experiments, but the graveyard entries in this benchmark --- particularly G5 (metaphor), G22 (gait normalization), and G26 (resurrected) --- are as methodologically informative as the successful findings. We recommend that future benchmark papers adopt a similar practice of documenting falsified hypotheses alongside confirmed ones.

\subsection{Connections to Theory}\label{connections-to-theory}

\textbf{\citet{karkada2025translation}.} They proved that static geometric structure in LLM representations arises inevitably from translation symmetry in co-occurrence statistics --- the manifold shape is determined by the data distribution, not by architectural choices. Our benchmark tests the dynamic complement: whether that geometry supports structured, consistent behavioral navigation. Together, the contributions are complementary.

Their result explains why static geometry exists; ours shows that it supports navigation that is structured and asymmetric at the universal level, but model-specific in its dynamics. Geometry comes from data statistics (their contribution); navigational behavior on that geometry is universal in structure but model-specific in gait (ours). The scale effect from Section 8 is consistent with their framework: if geometric structure requires sufficient capacity to extract translation symmetries from training data, smaller models (Llama 8B) may have the structure but lack the resolution to navigate through the same landmarks as frontier models.

\textbf{Gardenfors operationalized.} \citeauthor{gardenfors2000}'s (\citeyear{gardenfors2000}) conceptual spaces theory posits that concepts live in navigable geometric spaces with measurable distance. The benchmark provides a large-scale computational probe of this theory across 12 models and 100+ concept pairs. The partial confirmation is substantive: concepts do live in navigable spaces with measurable, compositional distance (triangle inequality approximately 91\%, hierarchical transitivity 4.9x over random). Polysemy sense differentiation (cross-pair Jaccard near zero for different senses of ``bank,'' ``bat,'' and ``crane'') confirms that paths reflect conceptual structure rather than surface word statistics. But the pervasive asymmetry (mean 0.811) challenges the standard metric assumption in conceptual spaces theory, which typically presumes symmetric distance functions. The data suggest that quasimetric structure --- metric-like but inherently directional --- may be a more appropriate formal framework for computational conceptual spaces.

\textbf{Platonic Representation Hypothesis refined.} \citet{huh2024platonic} predicted that larger models converge toward shared representations. Our behavioral test partially confirms this: navigational structure converges (all 12 models have gait and asymmetry), navigational landmarks converge for frontier-scale models (aggregate bridge frequency difference CI includes zero), but navigational gaits remain model-specific (the 2.9x gait range does not collapse with scale). The result establishes a boundary for the hypothesis: representational convergence does not imply behavioral convergence. Structure converges; gaits diversify. Partial confirmation with a clear boundary condition.

\textbf{\citet{wyss2025spectral}.} Spectral analysis of semantic attractors predicts discrete basins in LLM conceptual space. Our bridge bottlenecks and the ``deep basins'' observed in the word convergence game that inspired this benchmark \citep{hodge2025convergence} are plausibly analogous to those spectral attractors --- the behavioral manifestation of the basin structure their theory predicts. Bridge concepts like ``spectrum'' and ``dog'' behave as attractor basins --- navigational points that paths converge toward regardless of starting trajectory. The pre-fill displacement results (Section 6.1) demonstrate that these basins are causally active: they anchor paths when present and paths reorganize when they are displaced. The bridge frequency robustness finding (\textgreater{} 0.97 across all protocol conditions, O31) further supports the attractor interpretation --- these are deep structural features of the navigational landscape, not artifacts of a particular elicitation protocol.

\subsection{Future Work}\label{future-work}

The benchmark's findings and limitations suggest ten specific directions, ordered from the most directly motivated by current data to the most speculative.

\begin{enumerate}
\def\labelenumi{\arabic{enumi}.}
\item
  \textbf{Multi-variable bridge fragility model.} The mechanism ceiling (Section 7) was reached by testing single variables in isolation. A joint model --- bridge role x pre-fill content x dominance ratio x competitor count --- could capture the interaction effects that single-variable designs miss. The bimodal survival pattern (Section 6.1) suggests that at least two factors interact: ``sadness'' survives because it has no competitors and sits on an obligatory navigational axis, while ``harmony'' collapses because it competes with rhythm, pattern, and frequency in a crowded navigational space. The combinatorial cost is high (likely 5,000+ runs for a 4-factor design), but this is the only known path past the mechanism ceiling.
\item
  \textbf{Embedding-based distance validation.} Replace Jaccard with embedding cosine similarity to rescue the lexical limitation (Section 10.3, limitation 3). This would also test whether the cross-model distance failure (O18, O19) is a property of lexical metrics or of the navigational distances themselves.
\item
  \textbf{Multilingual conceptual topology.} Test whether the navigational geometry holds for non-English prompts and multilingual models. If the structure arises from data statistics \citep{karkada2025translation}, language-specific training distributions may produce different geometries.
\item
  \textbf{Scale threshold investigation.} The Llama family provides a natural within-family comparison: 8B (bridge frequency 0.200) vs Maverick at frontier scale (0.724). Intermediate checkpoints (Llama 70B, Llama 405B) would locate the scale threshold at which bridge content converges with frontier models and determine whether the transition is gradual or abrupt.
\item
  \textbf{Temporal stability under model updates.} The no-drift finding (O9) covers the data collection window. Longitudinal monitoring --- re-running the core pair set on the same model IDs at 3-month intervals --- would test whether navigational structure is preserved across training updates and provider-side model swaps.
\item
  \textbf{Human waypoint comparison.} A calibration subset of 10-15 concept pairs administered to human participants under the same prompt format would establish whether the structural properties (asymmetry, bridge bottlenecks, gait variation) are shared between human and LLM navigation or are specific to autoregressive generation. Given \citeauthor{tversky1977features}'s (\citeyear{tversky1977features}) finding that human similarity is asymmetric and \citeauthor{constantinescu2016grid}'s (\citeyear{constantinescu2016grid}) evidence for grid-cell-based conceptual navigation, partial overlap is plausible but untested.
\item
  \textbf{External anchoring.} Validate bridge predictions against ConceptNet relation strength, Small World of Words association norms, or Wikipedia link structure. This would partially address the self-grounding circularity (limitation 2) by grounding navigational landmarks in independently measured conceptual structure.
\item
  \textbf{Control pair redesign.} The single-pair strict-threshold approach has failed (limitation 1). Three alternatives: (a) relative consistency metrics comparing experimental vs control pairs on predictability rather than presence of structure; (b) pair batteries testing 20+ control pairs to characterize the distribution of spurious bridge frequencies; (c) difficulty-graded controls sampling from the full spectrum of conceptual distance.
\item
  \textbf{Chain-of-thought as intervention.} Eliciting waypoints with and without chain-of-thought reasoning would test whether explicit reasoning changes navigational structure or merely rationalizes the same implicit routes. This connects to the compositionality gap literature: if chain-of-thought improves multi-hop reasoning, it may also alter bridge selection, path composition, and the gait spectrum.
\item
  \textbf{Curvature estimation.} Within-model curvature is estimable from the triangle inequality data --- the distribution of triangle excess (d(A,B) + d(B,C) - d(A,C)) characterizes local geometry, with positive excess suggesting positive curvature and near-zero excess suggesting flatness. Cross-model curvature comparison is blocked by the distance failure (O18, O19) and should not be attempted until embedding-based distances are validated (item 2 above).
\end{enumerate}

\begin{center}\rule{0.5\linewidth}{0.5pt}\end{center}

The discussion synthesizes a benchmark that achieves high replication accuracy (approximately 80\%) for structural characterization while reaching a mechanism ceiling (approximately 15-20\%) for quantitative explanation. The core empirical contribution is a set of structural properties --- characteristic gaits spanning 2.9x across 12 models, quasimetric asymmetry at 0.811, bottleneck bridge topology with frequencies exceeding 0.97 under robustness testing, and compositional hierarchy with 4.9x transitivity over random controls --- that generalize across architectures and survive protocol variation. The core negative contribution is equally important: 7 falsified mechanistic hypotheses and 29 documented dead ends establish that the qualitative geometry of conceptual navigation is accessible to current methods while the quantitative mechanism is not. The benchmark's most serious limitation --- the control pair failure, with 0 of 24 replacement cells passing --- is acknowledged directly, and its validity is reframed as resting on predictability gaps rather than absolute baselines. The six theoretical connections --- to quasimetric topology, Tversky's asymmetry, Karkada et al.'s geometric inevitability, Gardenfors's conceptual spaces, the Platonic Representation Hypothesis, and spectral semantic attractors --- position the findings within a broader landscape where static geometry is well-characterized but dynamic navigation is only beginning to be understood. The most productive future directions involve multi-variable mechanism models, embedding-based distance validation, and human comparison --- each addressing a specific limitation identified across the six acts.
